Nas seis leis simuladas, a FWER ajustada da família original varia
de 1,6 a 3,1\%; na família temporal nova,
de 1,8 a 5,6\% (Tabela~\ref{tab:g-calibracao}).
No estresse AR0,6, a família temporal atinge 5,6\%,
Wilson 95\%[4,34;7,20], sem evidência de controle uniforme a 5\%.
São frequências sob nulidade global e desenho/bancos fixados, não
uma prova de controle forte sob quaisquer alternativas. Os limites
Wilson e erros Monte Carlo acompanham cada taxa; não se certifica
controle apenas por uma frequência pontual inferior a 5\%.

Os intervalos básicos falham de forma material. Entre os 84 casos
alvo--mecanismo, a cobertura marginal vai de 51,8
a 95,8\%. Por exemplo, com t5 heterocedástico eAR0,3,
a cobertura da média inicial é60,7\%, contra 95\% nominal.
Os intervalos são, portanto, diagnósticos aproximados, sem cobertura
nominal validada; não se usam seus extremos para delimitar efeitos
econômicos com confiança 95\%. Valores-p bilaterais centrados e ICs
básicos assimétricos não são procedimentos duais; seu desempenho
deve ser avaliado separadamente. A frequência ajustada simulada não
repara a falha de cobertura ou valida outras leis de erros.
Não se escolheu outro bloco ou estimador após observar essa falha.

No banco observado da média da década, a média bootstrap é $-0,119998$,
contra $-0,130668$ observado: deslocamento $+0,010670$, ou 8,17\%
da magnitude pontual. Esse deslocamento do banco não é um erro
financeiro identificado nem igual ao ponto médio dos quantis do IC.
